%! AP Statistics — Unit 1, Lesson 1.2 — Guided Notes
\documentclass[10pt]{article}
\usepackage{apstats-article}
\usepackage{apstats-boxes}

%=============================================================================
\begin{document}

\pageheader{Unit 1, Lesson 1.2}{Guided Notes}
\namedateperiod

\vspace{0.12in}

% ── OBJECTIVE ────────────────────────────────────────────────────────────────
\begin{objectivebox}
\small By the end of this lesson, I will be able to\dots\\[4pt]
\writeline\\[1pt]
\writeline\\[1pt]
\writeline
\end{objectivebox}

\vspace{0.1in}

% ── VOCABULARY ───────────────────────────────────────────────────────────────
\begin{vocabbox}
\begin{multicols}{2}
\termblanklong{Variable}
\termblanklong{Categorical variable}
\termblanklong{Quantitative variable}
\columnbreak
\termblanklong{Discrete quantitative variable}
\termblanklong{Continuous quantitative variable}
\termblanklong{Distribution of a variable}
\end{multicols}
\end{vocabbox}

\vspace{0.1in}

% ── HOOK ─────────────────────────────────────────────────────────────────────
\begin{hookbox}
\small\textbf{Student Data Table --- sort these 5 variables into two groups.
Do NOT look at the group names yet.}

\vspace{8pt}
\renewcommand{\arraystretch}{1.5}
\begin{tabularx}{\linewidth}{@{} C{0.18\linewidth} C{0.14\linewidth}
                                    C{0.14\linewidth} C{0.2\linewidth}
                                    C{0.2\linewidth} @{}}
\rowcolor{sky}
\textbf{Name} & \textbf{Grade} & \textbf{GPA} &
\textbf{Fav.\ Color} & \textbf{Sleep (hrs)} \\
\midrule
Jordan  & 11 & 3.7 & Blue   & 7.5 \\
Priya   & 10 & 3.2 & Green  & 6.0 \\
Marcus  & 12 & 2.9 & Red    & 8.0 \\
Sofia   & 11 & 3.5 & Purple & 7.0 \\
\end{tabularx}

\vspace{10pt}
\textbf{My two groups:}

\vspace{6pt}
\renewcommand{\arraystretch}{1.8}
\begin{tabularx}{\linewidth}{@{} |>{\centering\arraybackslash}X
                                    |>{\centering\arraybackslash}X| @{}}
\hline
\textbf{Group A} & \textbf{Group B} \\
\hline
\blank{5.5cm} & \blank{5.5cm} \\
\hline
\end{tabularx}

\vspace{8pt}
\textbf{What was your sorting rule?}\quad \blank{9cm}
\end{hookbox}

\vspace{0.1in}

% ── DIRECT INSTRUCTION ───────────────────────────────────────────────────────
\begin{notesbox}{Classifying Variables}
\small
\begin{multicols}{2}

\textbf{\textcolor{navy}{The Two Main Types}}\\[4pt]
\textbf{Categorical variable:}\\
Places an individual into one of several \blank{2.5cm}.
No meaningful arithmetic.\\[6pt]
\textit{Examples:}
\begin{itemize}[itemsep=2pt]
  \item Favorite color: \blank{2.2cm}
  \item Mode of transport: \blank{2.2cm}
\end{itemize}

\vspace{8pt}
\textbf{Quantitative variable:}\\
Takes numerical values for which \blank{2.5cm} makes sense.\\[6pt]
\textit{Examples:}
\begin{itemize}[itemsep=2pt]
  \item GPA: \blank{2.2cm}
  \item Hours of sleep: \blank{2.2cm}
\end{itemize}

\columnbreak

\textbf{\textcolor{navy}{The Average Test}}\\[4pt]
\textit{``Does averaging these values give a \blank{2.5cm} result?''}\\[6pt]
\begin{tabularx}{\linewidth}{@{} l X @{}}
\textbf{Yes} $\to$ & \blank{3cm} \\[4pt]
\textbf{No}  $\to$ & \blank{3cm} \\
\end{tabularx}

\vspace{10pt}
\textbf{\textcolor{navy}{Trap Variables --- Numbers That Are Categorical}}\\[4pt]
\begin{itemize}[itemsep=4pt]
  \item ZIP code: \blank{3cm}
  \item Jersey number: \blank{3cm}
  \item Student ID: \blank{3cm}
  \item Likert rating (1--5): \blank{2cm}
\end{itemize}
\end{multicols}
\end{notesbox}

\vspace{0.1in}

% ── DISCRETE VS CONTINUOUS ───────────────────────────────────────────────────
\begin{notesbox}{Discrete vs.\ Continuous Quantitative Variables}
\small
\begin{multicols}{2}

\textbf{\textcolor{navy}{Discrete}}\\
Values are \blank{2.5cm} --- there are \blank{2cm} between possible values.\\[4pt]
Ask: \textit{``Can I \blank{2.5cm} them?''}\\[6pt]
\textit{Example:} Number of pets\\
Possible values: \blank{3.5cm}\\[4pt]
\begin{itemize}[itemsep=2pt]
  \item Number of siblings: \blank{2cm}
  \item Number of songs on playlist: \blank{2cm}
\end{itemize}

\columnbreak

\textbf{\textcolor{navy}{Continuous}}\\
Can take \blank{2.5cm} value in an interval --- \blank{2.2cm}, not counted.\\[4pt]
Ask: \textit{``Can it fall \blank{2.5cm} two values?''}\\[6pt]
\textit{Example:} Weight of a pet\\
Possible values: \blank{3.5cm}\\[4pt]
\begin{itemize}[itemsep=2pt]
  \item Height in cm: \blank{2cm}
  \item Commute time: \blank{2cm}
\end{itemize}
\end{multicols}

\vspace{6pt}
\renewcommand{\arraystretch}{1.5}
\begin{tabularx}{\linewidth}{@{} >{\centering\arraybackslash\columncolor{navy}}p{0.28\linewidth}
                                    >{\centering\arraybackslash}X
                                    >{\centering\arraybackslash}X @{}}
\textcolor{white}{\textbf{Question}} &
\textcolor{navy}{\textbf{Discrete}} &
\textcolor{navy}{\textbf{Continuous}} \\
Can I list all possible values? & \blank{2cm} & \blank{2cm} \\
Is it measured on a scale?      & \blank{2cm} & \blank{2cm} \\
Are there gaps between values?  & \blank{2cm} & \blank{2cm} \\
\end{tabularx}
\end{notesbox}

\vspace{0.1in}

% ── DECISION FLOWCHART ───────────────────────────────────────────────────────
\begin{notesbox}{Decision Flowchart --- Copy \& Memorize}
\small\centering
\vspace{4pt}
\begin{tabular}{@{} c @{}}
\colorbox{navy}{\color{white}\textbf{\small Is the variable a word, label, or category?}}\\[4pt]
\end{tabular}

\vspace{2pt}
\begin{tabular}{@{} p{0.42\linewidth} @{\hspace{0.1\linewidth}} p{0.42\linewidth} @{}}
\centering\textbf{YES} $\Downarrow$ & \centering\textbf{NO} $\Downarrow$ \tabularnewline
\centering\colorbox{greenbg}{\textbf{Categorical}} &
\centering\textit{Does averaging give a meaningful result?} \tabularnewline
& \\
& \begin{tabular}{@{} p{0.45\linewidth} @{\hspace{0.05\linewidth}} p{0.45\linewidth} @{}}
  \textbf{YES} $\Downarrow$ & \textbf{NO} $\Downarrow$ \tabularnewline
  \textit{Can you count the values (are there gaps)?} & \colorbox{greenbg}{\textbf{Categorical}} \tabularnewline
  & \tabularnewline
  \end{tabular} \\
\end{tabular}

\vspace{4pt}
\begin{tabular}{@{} p{0.42\linewidth} @{\hspace{0.1\linewidth}} p{0.42\linewidth} @{}}
& \begin{tabular}{@{} p{0.45\linewidth} @{\hspace{0.05\linewidth}} p{0.45\linewidth} @{}}
  \textbf{YES} $\Downarrow$ & \textbf{NO} $\Downarrow$ \tabularnewline
  \colorbox{redbg}{\textbf{Quant.\ Discrete}} & \colorbox{redbg}{\textbf{Quant.\ Continuous}} \tabularnewline
  \end{tabular}
\end{tabular}
\end{notesbox}

\vspace{0.1in}

% ── GUIDED PRACTICE ──────────────────────────────────────────────────────────
\begin{practicebox}
\small\textbf{Classify each variable.} Write C, QD, or QC and give a one-phrase justification.

\vspace{6pt}
\renewcommand{\arraystretch}{2.0}
\begin{tabularx}{\linewidth}{@{} c >{\raggedright\arraybackslash}p{0.38\linewidth}
                                    C{0.1\linewidth} X @{}}
\rowcolor{sky}
\# & \textbf{Variable} & \textbf{Type} & \textbf{Justification} \\
\midrule
1 & Eye color              & \blank{1.2cm} & \blank{6.5cm} \\
2 & Number of siblings     & \blank{1.2cm} & \blank{6.5cm} \\
3 & Body temperature (°F)  & \blank{1.2cm} & \blank{6.5cm} \\
4 & Area code              & \blank{1.2cm} & \blank{6.5cm} \\
5 & Number of words in an essay & \blank{1.2cm} & \blank{6.5cm} \\
6 & Time to run a mile (min) & \blank{1.2cm} & \blank{6.5cm} \\
\end{tabularx}
\end{practicebox}

\vspace{0.1in}

% ── SPIRAL / BIG IDEAS ────────────────────────────────────────────────────────
\begin{spiralbox}
\small
\begin{multicols}{2}

\textbf{This lesson connects forward to\dots}
\begin{itemize}[itemsep=2pt]
  \item \textbf{Lessons 1.3--1.4:} Displaying categorical data (bar charts, two-way tables)
  \item \textbf{Lessons 1.5--1.9:} Displaying and summarizing quantitative data
  \item \textbf{Unit 4:} Discrete probability distributions vs.\ density curves
\end{itemize}

\columnbreak

\textbf{Big Idea --- Write in your own words:}\\
\textit{Why does variable type matter before we choose how to analyze data?}\\[2pt]
\writeline\\[1pt]\writeline\\[1pt]\writeline

\end{multicols}
\end{spiralbox}

\vspace{0.1in}

\begin{notesbox}{Additional Notes}
\vspace{0pt}
\writeline\\\writeline\\\writeline\\\writeline\\\writeline\\\writeline
\end{notesbox}

\end{document}
